\section{Discussion}
\label{sec:discussion}

\subsection{Interpreting the Null}

The reduced-form evidence does not support the hypothesis that user-fee-funded review acceleration changed the controlled-substance composition of FDA approvals. The aggregate CS share fell after PDUFA. The NDA-only Poisson rate model returns an incidence-rate ratio of $1.46$ ($p = 0.06$) under year fixed effects, but the same coefficient collapses to $1.19$ with a 95\% confidence interval of $[0.59, 2.39]$ once group-specific linear trends are allowed. The GDUFA result follows the same pattern in the opposite direction: a year-FE IRR of $0.668$ that flips to $1.38$ with confidence interval $[0.86, 2.21]$ under group-specific trends. In both cases the most flexible specification places the user-fee composition effect in a band that includes both small reductions and small increases relative to the pre-policy baseline. The point estimates flip sign across specifications --- a feature, not a bug, of the data: the year-FE estimates are absorbing slow-moving differences in the CS and non-CS approval streams that have no plausible structural relationship to the user-fee policy enacted at the boundary year.

This is the right way to read the result. It is not that the data are too noisy to detect anything. It is that, once the empirical specification is allowed to accommodate the secular trends that pharmaceutical innovation followed during the 1980s and 2010s, no level shift remains to be explained by user-fee acceleration. The negative finding therefore does what a precisely estimated null can do in policy economics: it bounds the joint effect that any plausible mechanism --- demand-side incentives operating through the present-value channel of Section~\ref{sec:theory}, supply-side capacity expansion through additional FDA staffing, or any combination of the two --- could have produced through the approval-composition margin.

\subsection{The 2009--2013 Opioid Product Cycle as a Separable Phenomenon}

The drug-level evidence in Table~\ref{tab:spike_drugs} makes a stronger case for a separable explanation than for a delayed PDUFA effect. Three features matter. First, the timing --- seventeen to twenty-one years after PDUFA enactment --- is inconsistent with an immediate or even moderately delayed regulatory-incentive response. Second, the drugs themselves are recognizable products from the documented opioid product-innovation cycle: reformulated OxyContin, Exalgo, Embeda, Suboxone, Zohydro ER, and a cluster of branded testosterone formulations. Third, the spike subsides after 2013, suggesting a product-cycle dynamic rather than a permanent shift in approval composition.

The spike inflates post-PDUFA IRRs in specifications that do not allow group-specific trends because the pre-PDUFA NDA CS rate was low and the late-2000s wave creates a large mean difference in the post period that the simpler models then attribute to the policy variable. The fact that the group-specific-trend specification eliminates this elevation confirms the diagnosis: the late-2000s wave is doing the work that the post-PDUFA dummy was getting credit for.

This finding fits the broader opioid-supply literature. \textcite{alpertSupplySideDrugPolicy2018a} document that the branded opioid wave was followed by a substitution response in illicit markets; the approval-side documentation here provides complementary upstream evidence. The implication for policy is that the regulatory approval channel and the industry innovation channel are distinct: FDA approval decisions necessarily preceded each opioid wave product, but the timing and clustering of those products are not naturally explained by user-fee acceleration. The driver appears to be firm-level investment decisions about a specific high-margin therapeutic segment --- not the PDUFA review-time incentive structure operating uniformly across all NDA applications.

\subsection{Interpreting the GDUFA Result}

The year-FE GDUFA estimate, taken on its own, would imply a one-third decline in the CS share of ANDA approvals after the post-2015 performance-goal era. The direction is opposite to the regulatory-acceleration hypothesis: if faster review disproportionately benefited CS generics, the CS ANDA share should rise, not fall. But the headline year-FE result is not robust. The group-specific-trends specification yields an IRR of $1.38$ with a confidence interval of $[0.86, 2.21]$, which is statistically indistinguishable from one. The level decline observed under year fixed effects is absorbed once the analysis accommodates the long pre-period decline in the CS share of ANDA approvals --- a decline that is visible in the descriptive series in Section~\ref{sec:results_gdufa} and that began well before GDUFA.

Three substantive forces are consistent with this pre-existing decline. First, the pre-GDUFA ANDA backlog disproportionately included CS generics filed during the 1990s--2000s opioid expansion; clearing the backlog during 2013--2014 produced elevated CS ANDA approvals in the transition period and a subsequent reversion as the backlog drained. Second, the post-2012 ANDA filing pipeline shifted toward non-CS generics in response to the 2005--2012 patent cliff for non-opioid branded drugs, mechanically lowering the CS share of approvals as those generics cleared review. Third, the prescription opioid market contracted after 2012 as prescribing guidelines tightened, state prescription drug monitoring programs expanded, and supply-side interventions restricted access. Reduced demand for opioid generics may have discouraged new CS ANDA filings relative to non-CS filings independent of GDUFA. These three channels are observationally close to one another and to the user-fee channel, and the data here cannot separate them.

The right interpretation of the GDUFA result is therefore the same as the PDUFA result: under the most flexible specification, the user-fee composition effect is indistinguishable from zero, and the direction of the year-FE point estimate provides little additional information once the slower-moving secular movement in the CS-ANDA share is allowed to load on its own trend.

\subsection{Limitations and Future Work}

Several limitations constrain the conclusions. The most important is the DEA linkage: because current schedules are used rather than historical schedules, drugs that were scheduled or rescheduled after their FDA approval date are misclassified. The conservative-match definition concentrates the linkage on substances that have been continuously scheduled for decades, which limits the practical extent of the bias, but the direction of any residual error for the main results is unclear. A second limitation is the single national time series, which provides limited statistical power and cannot separate the user-fee channel from contemporaneous secular trends in pharmaceutical innovation. A third limitation is the absence of a downstream outcome: this paper characterizes the approval-side composition of CS drugs but does not measure downstream prescribing volume, misuse, or diversion.

Four directions follow naturally for future work. First, constructing a historical DEA scheduling database would allow each drug to be classified by its schedule as of its approval date, eliminating the retroactive-linkage limitation. Second, linking approval-side CS counts to prescribing data from sources such as IQVIA or ARCOS would estimate the downstream supply implications of any composition change, even if that change is small. Third, extending the GDUFA analysis to GDUFA II and III as those program years accumulate will provide additional post-treatment observations and may sharpen the identification. Fourth, exploiting therapeutic-class or firm-level variation could begin to separate the demand-side incentive channel from the supply-side capacity channel within the user-fee mechanism, distinguishing the two competing theoretical accounts in Section~\ref{sec:theory}.
